% Source: Anteproyecto (3).pdf. Original wording retained except population/sample.
\chapter{Estado del arte}

 La revisión del estado del arte se organizó en tres ejes temáticos: (a) enfoques de aprendizaje autosupervisado aplicados a la clasificación de lesiones cutáneas, (b) sistemas de diagnóstico asistido por computador basados en aprendizaje profundo supervisado y transfer learning, y (c) investigaciones desarrolladas en el contexto colombiano sobre inteligencia artificial y dermatología. Esta estructura permite contrastar las aproximaciones existentes con la propuesta del presente proyecto e identificar los vacíos que justifican su desarrollo.

\section{Aprendizaje autosupervisado en la clasificación de lesiones cutáneas}

 El aprendizaje autosupervisado (SSL) ha emergido como una alternativa prometedora para reducir la dependencia de datos etiquetados en el análisis de imágenes dermatológicas. \textcite{useini2024automatizedSSL} propusieron un sistema automatizado de screening basado en DINO con backbone ResNet18, capaz de detectar y caracterizar lesiones sospechosas tipo "ugly duckling" mediante clustering autosupervisado, sin requerir etiquetas manuales. Aunque los resultados demostraron la viabilidad del SSL para la detección sin supervisión, el estudio se limitó a pacientes de alto riesgo con más de 100 nevos y no fue validado en poblaciones con fototipos diversos (IV-V), condición prevalente en Colombia.

En la misma línea, \textcite{harczos2024barlowtwinsisic} compararon directamente el preentrenamiento autosupervisado mediante Barlow Twins (ResNet-50) frente al transfer learning supervisado desde ImageNet, utilizando el dataset ISIC 2019 con 25,331 imágenes. Los resultados evidenciaron que el modelo SSL alcanzó una precisión media de 0.703, superando al supervisado (0.663) en escenarios de datos etiquetados limitados. Este hallazgo es particularmente relevante para el contexto de Santander, donde la disponibilidad de anotaciones clínicas es restringida. No obstante, los autores reconocen que el rendimiento aún no es suficiente para aplicaciones clínicas directas.

El avance más significativo en esta área lo representa PanDerm, un modelo fundacional multimodal desarrollado por \textcite{yan2025panderm} y publicado en Nature Medicine. PanDerm fue preentrenado mediante masked latent modeling y alineación CLIP sobre más de 2 millones de imágenes de 11 instituciones clínicas, abarcando cuatro modalidades (fotografía de cuerpo completo, dermoscopia, clínica y dermatopatología). El modelo alcanzó el estado del arte en 28 benchmarks diversos, superando a los clínicos en un 10.2\% en la detección temprana de melanoma y mejorando el diagnóstico de profesionales no dermatólogos en un 16.5\%, incluso utilizando solo el 10\% de los datos etiquetados. Esta investigación confirma que el SSL puede superar al aprendizaje supervisado completo cuando se combina con fine-tuning eficiente, aunque su replicabilidad se ve limitada por los recursos computacionales masivos requeridos y el acceso restringido al dataset propietario.

\section{Sistemas CAD basados en aprendizaje profundo supervisado y transfer learning}

 En el ámbito del aprendizaje supervisado, los modelos basados en transfer learning han consolidado su posición como referentes en la clasificación de lesiones cutáneas. Un estudio reciente integró EfficientNet-B0, EfficientNet-B2 y ResNet50 mediante un ensemble con transfer learning desde ImageNet, alcanzando una accuracy máxima de 99.14\% en datasets combinados de HAM10000, ISIC y Kaggle (Diagnostics, 2025 % TODO citas: DOI confirmado (10.3390/diagnostics15050551, ver fila de la Tabla~\ref{tab:estado-arte}), autor y titulo no confirmados; falta agregar entrada real a referencias.bib
). Individualmente, EfficientNet-B2 alcanzó un 89.55\% de accuracy. Estos resultados, aunque notables, dependen de conjuntos de datos completamente etiquetados, lo que los hace difícilmente replicables en contextos donde la anotación clínica es costosa y limitada, como en el departamento de Santander. Complementariamente, \textcite{himel2024skintrans} implementaron Vision Transformers (ViT) combinados con SkinTrans sobre el dataset HAM10000, reportando una precisión del 94.3\% en la clasificación de siete tipos de lesiones. Este trabajo evidencia el potencial de las arquitecturas basadas en atención para captar relaciones espaciales globales en las imágenes dermatológicas, aunque mantiene la limitación de requerir datos completamente supervisados. Estos estudios establecen el benchmark supervisado contra el cual debe compararse el modelo autosupervisado propuesto en este proyecto, utilizando las mismas familias de arquitecturas (ResNet, EfficientNet, ViT) como backbone.

\section{Investigaciones en el contexto colombiano}

 En el ámbito nacional, se identificaron contribuciones relevantes aunque aún incipientes. \textcite{riosduarte2024colombianmelanoma}, de la Universidad de los Andes (Bogotá), realizaron un análisis comparativo de modelos CNN entrenados con imágenes clínicas, dermatoscópicas y combinadas para la clasificación de melanoma. Sus resultados, publicados en Skin Research and Technology, demostraron que los modelos entrenados con dermatoscopía superan a los basados únicamente en imágenes clínicas, mientras que la combinación de ambas modalidades no evidenció ventajas claras. En resumen, este trabajo propone un método de clasificación supervisada de melanoma con CNN en el contexto colombiano.

\textcite{florez2022clasificaciondermatologica}, en la Revista Colombiana de Tecnologías de Avanzada, implementaron cuatro métodos de machine learning clásico (SVM, KNN, redes neuronales y árboles de decisión) para la clasificación de imágenes del dataset HAM10000. Si bien reportaron precisiones cercanas al 100\% en muestras balanceadas, los resultados presentaron alta variabilidad, evidenciando las limitaciones inherentes al aprendizaje automático tradicional frente a técnicas de aprendizaje profundo y, particularmente, frente al aprendizaje autosupervisado.

\textcite{torresospina2025iaprisma} publicaron en la revista Ibero Ciencias una revisión sistemática basada en la metodología PRISMA, analizando 25 estudios sobre la integración de la inteligencia artificial en la detección temprana del cáncer de piel. Los autores concluyeron que las CNN alcanzan niveles de precisión comparables o superiores a los dermatólogos, pero identificaron desafíos persistentes en validación clínica, sesgo algorítmico y aceptación ética. Notablemente, ninguno de los estudios revisados empleó técnicas de aprendizaje autosupervisado, lo que confirma el vacío que el presente proyecto busca abordar.

Finalmente, una revisión panorámica publicada en \textcite{biomedica2025iasaludcolombia}, la revista del Instituto Nacional de Salud de Colombia, documentó el estado actual del ecosistema de inteligencia artificial y salud digital en el país. El análisis evidenció que, aunque existen avances significativos en investigación aplicada, la transición de experiencias piloto a implementación clínica efectiva enfrenta múltiples barreras, incluyendo limitaciones regulatorias, de infraestructura y de interoperabilidad. Este panorama refuerza la pertinencia de desarrollar herramientas computacionales adaptadas al contexto regional colombiano.

\section{Síntesis}

La revisión del estado del arte, sintetizada en la Tabla~\ref{tab:estado-arte}, permite identificar tres vacíos que justifican el desarrollo del presente proyecto. En primer lugar, ninguna de las investigaciones nacionales identificadas emplea técnicas de aprendizaje autosupervisado, concentrándose exclusivamente en enfoques supervisados clásicos o métodos de machine learning tradicional, lo que evidencia una ausencia de aplicaciones SSL en el contexto colombiano. En segundo lugar, los estudios internacionales de SSL fueron desarrollados y evaluados predominantemente en poblaciones caucásicas, sin considerar los fototipos IV-V prevalentes en Colombia, lo que limita su extrapolación directa al contexto regional. En tercer lugar, aunque modelos como PanDerm demuestran resultados excepcionales, su despliegue requiere recursos computacionales fuera del alcance de instituciones clínicas regionales, evidenciando una brecha significativa entre el rendimiento teórico documentado en la literatura y las condiciones reales de implementación en entornos de recursos limitados.

El presente proyecto se posiciona en la intersección de estas brechas, proponiendo la adaptación e implementación de un modelo de aprendizaje autosupervisado entrenado con conjuntos de datos públicos, evaluado mediante métricas estándar de clasificación y comparado con una línea base supervisada, con el propósito de desarrollar un prototipo funcional de apoyo diagnóstico dermatológico orientado al contexto clínico del departamento de Santander.

\begingroup
\small
\singlespacing
\setlength{\tabcolsep}{4pt}
\begin{longtable}{@{}>{\raggedright\arraybackslash}p{0.14\textwidth}>{\raggedright\arraybackslash}p{0.18\textwidth}>{\raggedright\arraybackslash}p{0.17\textwidth}>{\raggedright\arraybackslash}p{0.13\textwidth}>{\raggedright\arraybackslash}p{0.28\textwidth}@{}}
\caption{Estado del arte del aprendizaje autosupervisado en dermatología y cáncer de piel}\label{tab:estado-arte}\\
\toprule
\textbf{Categoría} & \textbf{Referencia (Autores, Año, País, DOI)} & \textbf{Metodología / Arquitectura} & \textbf{Dataset} & \textbf{Resultados} \\
\midrule
\endfirsthead
\toprule
\textbf{Categoría} & \textbf{Referencia (Autores, Año, País, DOI)} & \textbf{Metodología / Arquitectura} & \textbf{Dataset} & \textbf{Resultados} \\
\midrule
\endhead
\bottomrule
\endfoot
SSL en Dermatología & Fabian, V. et al., 2024 (Suiza). \url{https://doi.org/10.1038/s41598-024-61681-4} & DINO + ResNet18; clustering autosupervisado & 90 pacientes (FotoFinder ATBM) & Detecta y caracteriza lesiones sin etiquetas mediante agrupación visual. Limitación: datos poco diversos (fototipos altos no evaluados). Aporte: valida el uso de SSL (DINO) sin anotaciones, alineado con el enfoque del proyecto. \\
\addlinespace
SSL con Datos Limitados & Harczos, F. et al., 2024 (Alemania). \url{https://arxiv.org/abs/2401.00692} & Barlow Twins + ResNet50 & ISIC 2019 & SSL alcanza mayor accuracy (0.703) que supervisado (0.663), demostrando mejor desempeño con pocos datos etiquetados y alta transferibilidad. Limitación: rendimiento aún no clínico y arquitectura limitada. Aporte: evidencia clave para escenarios con escasez de datos (como Santander). \\
\addlinespace
Modelo Fundacional SSL & Yan, S. et al., 2025 (Australia). \url{https://doi.org/10.1038/s41591-025-03747-y} & Masked latent + CLIP multimodal & 2M imágenes (11 instituciones) & Logra SOTA en 28 benchmarks, supera clínicos en melanoma temprano (10.2\%) y mejora no especialistas (16.5\%), mostrando el potencial de modelos fundacionales multimodales. Limitación: alto costo computacional y datos propietarios. Aporte: referencia de vanguardia que demuestra superioridad del SSL con pocos datos etiquetados. \\
\addlinespace
Deep Learning Colombia & Ríos-Duarte, J.A. et al., 2024 (Colombia). \url{https://doi.org/10.1111/srt.13607} & CNN + transfer learning & Imágenes clínicas y dermoscópicas & La dermatoscopía supera imágenes clínicas, y la combinación no mejora significativamente, destacando la importancia del tipo de imagen. Limitación: no usa SSL ni evalúa diversidad de fototipos. Aporte: línea base supervisada en Colombia para comparación futura con SSL. \\
\addlinespace
Machine Learning Colombia & Flórez Fuentes, A.S. et al., 2022 (Colombia). \url{https://ojs.unipamplona.edu.co/index.php/rcta/article/view/2350} & SVM, KNN, RN, árboles & HAM10000 & Modelos clásicos alcanzan alta precisión en datos balanceados, pero con gran variabilidad en resultados. Limitación: baja robustez y ausencia de deep learning/SSL. Aporte: evidencia la necesidad de métodos más sólidos como SSL. \\
\addlinespace
Revisión IA Colombia & Torres Ospina, M. et al., 2025 (Colombia). \url{https://revistaiberociencias.org/index.php/multidisciplinar/article/view/33} & Revisión PRISMA & N/A & CNNs logran precisión comparable a dermatólogos y crece el uso de apps diagnósticas. Limitación: problemas de validación clínica, sesgo y falta de SSL. Aporte: identifica vacío en SSL en dermatología en Colombia. \\
\addlinespace
IA en Salud Colombia & Biomédica, 2025 (Colombia). \url{https://revistabiomedica.org/index.php/biomedica/article/view/8248} & Revisión panorámica & N/A & El ecosistema de IA en salud está en consolidación, con avances en investigación pero barreras en implementación. Limitación: regulación, infraestructura e interoperabilidad. Aporte: contextualiza la necesidad de soluciones adaptadas al entorno colombiano. \\
\addlinespace
Transfer Learning Dermatología & Diagnostics, 2025 (Internacional). \url{https://doi.org/10.3390/diagnostics15050551} & EfficientNet + ResNet (ensemble) & HAM10000, ISIC & Logra hasta 99.14\% de accuracy con modelos fusionados, evidenciando el alto rendimiento del enfoque supervisado. Limitación: requiere datos completamente etiquetados. Aporte: benchmark para comparar futuros modelos SSL del proyecto. \\
\addlinespace
Deep Learning Colombia & Calderón, C.; Sánchez, K.; Argüello, H., 2021 (Colombia). \url{https://doi.org/10.1016/j.cmpbup.2021.100036} & CNN bilineal: ResNet50 + VGG16 + Xception; transfer learning con fine-tuning & HAM10000 (7 clases) & Accuracy de 0.9321 (mejora 2.7\% sobre SOTA previo), AUC de 0.9810 mediante fusión bilineal de CNNs con data augmentation. Limitación: solo supervisado, no evalúa fototipos diversos. Aporte: co-autoría de la asesora del proyecto (Karen Sánchez); benchmark supervisado con mismo dataset HAM10000 que será usado en el proyecto SSL. \\
\addlinespace
\end{longtable}
\endgroup
   Nota. Elaboración propia a partir de fuentes científicas \parencite{useini2024automatizedSSL,harczos2024barlowtwinsisic,yan2025panderm,riosduarte2024colombianmelanoma,florez2022clasificaciondermatologica,torresospina2025iaprisma,biomedica2025iasaludcolombia} y la fuente de Diagnostics (2025) señalada como pendiente de confirmación bibliográfica más arriba en este capítulo.
